% Supplementary Material for "The Geometry of Consolidation"
% Authors: Anirudh Bharadwaj Vangara (Sentra; University of Waterloo Computer Engineering),
%          Ashwin Gopinath (Sentra; MIT Mechanical Engineering, corresponding).
%
% Compiled standalone. For one-PDF submission, append to main.tex.

\documentclass[11pt]{article}

\usepackage[margin=1.0in]{geometry}
\usepackage{amsmath,amssymb,amsthm,mathtools}
\usepackage{graphicx}
\usepackage{hyperref}
\usepackage{xcolor}
\usepackage{booktabs}
\usepackage{longtable}
% microtype disabled: enabling font expansion triggers pdftex fatal
% error on non-scalable bitmap fonts used by the long per-cell tables.
% \usepackage{microtype}
\usepackage{authblk}
\usepackage[sort&compress,numbers,super]{natbib}
\usepackage{titlesec}
\usepackage{enumitem}

\titleformat*{\section}{\normalsize\bfseries}
\titleformat*{\subsection}{\normalsize\itshape}

\newtheorem{theorem}{Theorem}
\newtheorem{lemma}[theorem]{Lemma}
\newtheorem{proposition}[theorem]{Proposition}
\newtheorem{corollary}[theorem]{Corollary}
\theoremstyle{definition}
\newtheorem{definition}[theorem]{Definition}
\theoremstyle{remark}
\newtheorem*{remark}{Remark}

\newcommand{\deff}{d_{\mathrm{eff}}}
\newcommand{\dbar}{\bar{d}}
\newcommand{\thetap}{\theta'}
\newcommand{\memset}{\mathcal{M}}
\newcommand{\clusterset}{\mathcal{C}}
\newcommand{\store}{\mathcal{S}}
\newcommand{\reps}{\mathcal{R}}
\newcommand{\norm}[1]{\left\lVert#1\right\rVert}
\newcommand{\inner}[2]{\left\langle#1,#2\right\rangle}
\newcommand{\eps}{\varepsilon}
\newcommand{\R}{\mathbb{R}}

\title{The Geometry of Consolidation \\
       \large Supplementary Material}

\author[1,2]{Anirudh Bharadwaj Vangara}
\author[1,3]{Ashwin Gopinath\thanks{\texttt{ashwin@sentra.app}}}
\affil[1]{Sentra, 235 2nd Street, San Francisco, CA 94105, USA}
\affil[2]{Department of Computer Engineering, University of Waterloo, Waterloo, ON N2L 3G1, Canada}
\affil[3]{Department of Mechanical Engineering, Massachusetts Institute of Technology, Cambridge, MA 02139, USA}

\date{Submitted to \textbf{NeurIPS 2026} as supplementary material.}

\begin{document}

\maketitle

% Re-number so supp. figures/tables do not collide with the main paper.
\renewcommand{\thefigure}{S\arabic{figure}}
\renewcommand{\thetable}{S\arabic{table}}
\renewcommand{\theequation}{S\arabic{equation}}
\setcounter{figure}{0}
\setcounter{table}{0}
\setcounter{equation}{0}
\setcounter{section}{0}
\renewcommand{\thesection}{S\arabic{section}}

\tableofcontents

\vspace{1em}

\section{Scope of this supplementary}

This document is a companion to the main paper; it does not restate
its narrative. It collects (i)~per-cell tables for every experimental
grid referenced in the main paper, (ii)~an expanded E8 encoder grid
with two additional figures, (iii)~the full $c_1$ calibration table at
every $(\deff, \theta)$ row of the E1 grid, (iv)~a run manifest mapping
every main-paper figure to its source parquet, and (v)~the exact
reproduction recipe. The proof of the Duality Theorem and the
statement of Definition~1 ($\deff$ as the participation ratio
$(\sum\lambda_i)^2/\sum\lambda_i^2$) live in the main paper's Methods
section and are not duplicated here.

\paragraph{What is \emph{not} in this supplementary.}
No new theorems. No additional experiments beyond those referenced in
the main paper. No claims about $c_1$ or the regime boundary that are
not already stated quantitatively in \S3 and \S4 of the main paper.

\section{Expanded E8 encoder-strategy grid}
\label{sec:supp:e8-expanded}

The main paper reports E8 on a 5-strategy slice ($6$ encoders $\times$
$5$ strategies $\times$ $3$ seeds $= 90$ cells). For completeness we
re-ran E8 on the full $6\times 11 \times 3 = 198$-cell grid that also
sweeps GAC's threshold $\theta \in \{0.6, 0.7, 0.8, 0.9\}$, four
\textsc{selective\_prune} keep-ratios ($0.50, 0.25, 0.10$), and the
\emph{no\_consolidation} identity ceiling. Two figures summarise the
expanded data.

\begin{figure}[h]
  \centering
  \includegraphics[width=0.78\linewidth]{figs/figS_theta_sweep.pdf}
  \caption{E8 expanded: GAC threshold $\theta$ sweep on DRM. Each line
  is one encoder; the dashed line is the mean centroid baseline.
  Identity accuracy is approximately flat across $\theta \in [0.6,0.9]$
  for every encoder --- consistent with the main paper's finding that
  tight DRM clusters fall in the regime where the centroid branch of
  the GAC router dominates, making the router insensitive to $\theta$.}
  \label{fig:s-theta}
\end{figure}

\begin{figure}[h]
  \centering
  \includegraphics[width=0.82\linewidth]{figs/figS_compression_frontier.pdf}
  \caption{E8 expanded: compression--identity frontier across the full
  $6 \times 11 \times 3$ grid. Each point is one seed. The horizontal
  axis is the empirical compression ratio
  $n_{\mathrm{train}}/n_{\mathrm{representatives}}$; the right edge is
  \emph{no\_consolidation} by construction. Centroid and medoid
  occupy the upper-right (compressed + accurate) corner on every
  encoder; \textsc{selective\_prune} degrades continuously as the
  keep-ratio shrinks.}
  \label{fig:s-frontier}
\end{figure}

\section{E1 strategy sweep and identity--coverage scatter}
\label{sec:supp:e1-aux}

Two auxiliary plots that did not fit in the main text.

\begin{figure}[h]
  \centering
  \includegraphics[width=0.92\linewidth]{figs/figS_strategy_sweep.pdf}
  \caption{Strategy-level cap-coverage error on the E1 grid. Each row
  fixes $\theta$; each panel plots mean $\eps_{\mathrm{cap}}$ versus
  $\deff$ for centroid, medoid, selective pruning, and GAC. All
  strategies collapse to zero error in the tight regime ($\dbar <
  \thetap$, left of the boundary in each panel); the four strategies
  fan out in the spread regime in the $\deff$-monotone order the
  Duality Theorem anticipates.}
  \label{fig:s-strategy-sweep}
\end{figure}

\begin{figure}[h]
  \centering
  \includegraphics[width=0.82\linewidth]{figs/figS_id_coverage.pdf}
  \caption{Identity accuracy versus cap-coverage on the E4 corpus
  grid. Each point is one (corpus, strategy, encoder) triple. Identity
  and cap-coverage move together in the tight regime (lower-left
  cluster, near-zero error on both) and separate in the spread regime
  --- a direct visualisation that cap-coverage and identity are
  distinct quantities but agree wherever the theorem is non-vacuous.}
  \label{fig:s-id-cov}
\end{figure}

\section{Complete $c_1$ calibration table}
\label{sec:supp:c1}

Table~\ref{tab:c1_calib} gives the per-cell $c_1$ calibration row by
row across the E1 grid. The aggregate numbers reported in \S4 of the
main paper ($c_1^{95} \approx 0.046$ on $12{,}400$ tight cells;
$c_1^{95} \approx 4.61 \times 10^{10}$ on $3{,}600$ spread cells;
global $c_1^{95} \approx 4.60 \times 10^{6}$ on all $16{,}000$ cells)
are the $p95$ quantiles of the per-row $c_1$ distribution restricted to
the relevant regime. The full per-$(d_{\mathrm{eff,target}}, \theta)$
breakdown below lets a reader audit the regime boundary and verify
that the $c_1$ explosion is localised to the extreme-coherence corner
(high $\theta$, low $\deff$), not an artifact of outlier rows.

\begin{table}[t]
\centering
\small
\caption{E1 calibration table: bound vs observed cap-coverage error across $(d_{eff}, \theta)$.}
\label{tab:c1_calib}
\begin{tabular}{lcccc}
\toprule
d\_eff\_target & theta & bound & observed\_err & theorem\_holds \\
\midrule
4.000 & 0.700 & 0.900 & 0.021 & 0.900 \\
4.000 & 0.800 & 0.850 & 0.044 & 0.850 \\
4.000 & 0.900 & 0.750 & 0.101 & 0.750 \\
4.000 & 0.950 & 0.637 & 0.170 & 0.637 \\
8.000 & 0.700 & 0.900 & 0.024 & 0.900 \\
8.000 & 0.800 & 0.850 & 0.055 & 0.850 \\
8.000 & 0.900 & 0.750 & 0.124 & 0.750 \\
8.000 & 0.950 & 0.619 & 0.196 & 0.619 \\
16.000 & 0.700 & 0.900 & 0.030 & 0.900 \\
16.000 & 0.800 & 0.850 & 0.069 & 0.850 \\
16.000 & 0.900 & 0.752 & 0.139 & 0.752 \\
16.000 & 0.950 & 0.635 & 0.216 & 0.635 \\
32.000 & 0.700 & 0.900 & 0.038 & 0.900 \\
32.000 & 0.800 & 0.850 & 0.076 & 0.850 \\
32.000 & 0.900 & 0.762 & 0.151 & 0.762 \\
32.000 & 0.950 & 0.635 & 0.228 & 0.635 \\
64.000 & 0.700 & 0.900 & 0.043 & 0.900 \\
64.000 & 0.800 & 0.850 & 0.083 & 0.850 \\
64.000 & 0.900 & 0.774 & 0.159 & 0.774 \\
64.000 & 0.950 & 0.644 & 0.237 & 0.644 \\
\bottomrule
\end{tabular}
\end{table}


\section{Per-cell experiment tables}
\label{sec:supp:percell}

These tables enumerate every row of the main experimental grids so
that aggregated statistics in the main paper can be reproduced by
inspection.

\begin{longtable}{lccccccccc}
\caption{E2 per-cell results (strategy sweep across 6 corpora, 6 encoders, 4 clusterers). (showing first 60 rows; full CSV in repo)}\label{tab:supp_e2}\\
\toprule
corpus & model & clusterer & strategy & seed & n\_train & compression & identity\_accuracy & recall@5 & coverage@0.80 \\
\midrule
\endfirsthead
\multicolumn{10}{l}{\small\itshape (continued)}\\
\toprule
corpus & model & clusterer & strategy & seed & n\_train & compression & identity\_accuracy & recall@5 & coverage@0.80 \\
\midrule
\endhead
\bottomrule
\endfoot
arxiv\_titles & bge-large & kmeans & centroid & 0 & 25545 & 2322.27 & 0.754 & 0.979 & 0.733 \\
arxiv\_titles & bge-large & hdbscan & centroid & 0 & 25545 & 2322.27 & 0.754 & 0.979 & 0.733 \\
arxiv\_titles & bge-large & hdbscan & gac & 0 & 25545 & 2.00 & 0.582 & 0.868 & 0.612 \\
arxiv\_titles & bge-large & kmeans & gac & 0 & 25545 & 2.00 & 0.582 & 0.868 & 0.612 \\
arxiv\_titles & bge-large & hdbscan & gac & 0 & 25545 & 2.00 & 0.582 & 0.868 & 0.612 \\
arxiv\_titles & bge-large & kmeans & gac & 0 & 25545 & 2.00 & 0.582 & 0.868 & 0.612 \\
arxiv\_titles & bge-large & kmeans & importance\_weighted & 0 & 25545 & 2322.27 & 0.754 & 0.980 & 0.734 \\
arxiv\_titles & bge-large & hdbscan & importance\_weighted & 0 & 25545 & 2322.27 & 0.754 & 0.980 & 0.734 \\
arxiv\_titles & bge-large & kmeans & medoid & 0 & 25545 & 2322.27 & 0.611 & 0.937 & 0.073 \\
arxiv\_titles & bge-large & hdbscan & medoid & 0 & 25545 & 2322.27 & 0.611 & 0.937 & 0.073 \\
arxiv\_titles & bge-large & kmeans & selective\_prune & 0 & 25545 & 2.00 & 0.582 & 0.868 & 0.612 \\
arxiv\_titles & bge-large & hdbscan & selective\_prune & 0 & 25545 & 2.00 & 0.582 & 0.868 & 0.612 \\
arxiv\_titles & bge-large & kmeans & selective\_prune & 0 & 25545 & 4.00 & 0.442 & 0.761 & 0.326 \\
arxiv\_titles & bge-large & hdbscan & selective\_prune & 0 & 25545 & 4.00 & 0.442 & 0.761 & 0.326 \\
c4\_random & bge-large & kmeans & centroid & 0 & 8787 & 17.57 & 0.713 & 0.948 & 0.082 \\
c4\_random & bge-large & kmeans & centroid & 0 & 8787 & 17.57 & 0.012 & 0.058 & 0.000 \\
c4\_random & bge-large & kmeans & centroid & 0 & 8787 & 17.57 & 0.002 & 0.018 & 0.000 \\
c4\_random & bge-large & kmeans & centroid & 0 & 8787 & 17.57 & 0.002 & 0.013 & 0.000 \\
c4\_random & bge-large & kmeans & gac & 0 & 8787 & 4.27 & 0.626 & 0.813 & 0.003 \\
c4\_random & bge-large & kmeans & gac & 0 & 8787 & 4.27 & 0.006 & 0.021 & 0.000 \\
c4\_random & bge-large & kmeans & gac & 0 & 8787 & 4.27 & 0.002 & 0.009 & 0.000 \\
c4\_random & bge-large & kmeans & gac & 0 & 8787 & 4.27 & 0.002 & 0.007 & 0.000 \\
c4\_random & bge-large & kmeans & gac & 0 & 8787 & 2.04 & 0.312 & 0.583 & 0.004 \\
c4\_random & bge-large & kmeans & gac & 0 & 8787 & 2.04 & 0.005 & 0.031 & 0.000 \\
c4\_random & bge-large & kmeans & gac & 0 & 8787 & 2.04 & 0.005 & 0.016 & 0.000 \\
c4\_random & bge-large & kmeans & gac & 0 & 8787 & 2.04 & 0.004 & 0.013 & 0.000 \\
c4\_random & bge-large & kmeans & medoid & 0 & 8787 & 17.57 & 0.650 & 0.881 & 0.002 \\
c4\_random & bge-large & kmeans & medoid & 0 & 8787 & 17.57 & 0.010 & 0.045 & 0.000 \\
c4\_random & bge-large & kmeans & medoid & 0 & 8787 & 17.57 & 0.003 & 0.018 & 0.000 \\
c4\_random & bge-large & kmeans & medoid & 0 & 8787 & 17.57 & 0.002 & 0.013 & 0.000 \\
c4\_random & bge-large & kmeans & selective\_prune & 0 & 8787 & 4.49 & 0.206 & 0.467 & 0.002 \\
c4\_random & bge-large & kmeans & selective\_prune & 0 & 8787 & 4.49 & 0.007 & 0.035 & 0.000 \\
c4\_random & bge-large & kmeans & selective\_prune & 0 & 8787 & 4.49 & 0.005 & 0.021 & 0.000 \\
c4\_random & bge-large & kmeans & selective\_prune & 0 & 8787 & 4.49 & 0.004 & 0.017 & 0.000 \\
c4\_random & bge-large & kmeans & selective\_prune & 0 & 8787 & 8.08 & 0.167 & 0.404 & 0.001 \\
c4\_random & bge-large & kmeans & selective\_prune & 0 & 8787 & 8.08 & 0.008 & 0.031 & 0.000 \\
c4\_random & bge-large & kmeans & selective\_prune & 0 & 8787 & 8.08 & 0.007 & 0.020 & 0.000 \\
c4\_random & bge-large & kmeans & selective\_prune & 0 & 8787 & 8.08 & 0.006 & 0.016 & 0.000 \\
drm\_templated & bge-large & hdbscan & centroid & 0 & 3200 & 4.00 & 0.943 & 0.996 & 1.000 \\
drm\_templated & bge-large & kmeans & centroid & 0 & 3200 & 4.00 & 0.943 & 0.996 & 1.000 \\
drm\_templated & bge-large & kmeans & gac & 0 & 3200 & 1.55 & 0.932 & 0.986 & 1.000 \\
drm\_templated & bge-large & hdbscan & gac & 0 & 3200 & 1.55 & 0.932 & 0.986 & 1.000 \\
drm\_templated & bge-large & kmeans & gac & 0 & 3200 & 1.34 & 0.932 & 0.984 & 1.000 \\
drm\_templated & bge-large & hdbscan & gac & 0 & 3200 & 1.34 & 0.932 & 0.984 & 1.000 \\
drm\_templated & bge-large & kmeans & importance\_weighted & 0 & 3200 & 4.00 & 0.941 & 0.996 & 1.000 \\
drm\_templated & bge-large & hdbscan & importance\_weighted & 0 & 3200 & 4.00 & 0.941 & 0.996 & 1.000 \\
drm\_templated & bge-large & kmeans & medoid & 0 & 3200 & 4.00 & 0.910 & 0.993 & 1.000 \\
drm\_templated & bge-large & hdbscan & medoid & 0 & 3200 & 4.00 & 0.910 & 0.993 & 1.000 \\
drm\_templated & bge-large & kmeans & selective\_prune & 0 & 3200 & 2.00 & 0.769 & 0.884 & 1.000 \\
drm\_templated & bge-large & hdbscan & selective\_prune & 0 & 3200 & 2.00 & 0.769 & 0.884 & 1.000 \\
drm\_templated & bge-large & kmeans & selective\_prune & 0 & 3200 & 4.00 & 0.761 & 0.886 & 1.000 \\
drm\_templated & bge-large & hdbscan & selective\_prune & 0 & 3200 & 4.00 & 0.761 & 0.886 & 1.000 \\
drm\_templated & e5-large & kmeans & centroid & 0 & 3200 & 4.00 & 0.965 & 0.999 & 1.000 \\
drm\_templated & e5-large & hdbscan & centroid & 0 & 3200 & 4.00 & 0.965 & 0.999 & 1.000 \\
drm\_templated & e5-large & kmeans & gac & 0 & 3200 & 1.95 & 0.963 & 0.999 & 1.000 \\
drm\_templated & e5-large & hdbscan & gac & 0 & 3200 & 1.95 & 0.963 & 0.999 & 1.000 \\
drm\_templated & e5-large & kmeans & gac & 0 & 3200 & 1.95 & 0.963 & 0.999 & 1.000 \\
drm\_templated & e5-large & hdbscan & gac & 0 & 3200 & 1.95 & 0.963 & 0.999 & 1.000 \\
drm\_templated & e5-large & hdbscan & importance\_weighted & 0 & 3200 & 4.00 & 0.965 & 0.999 & 1.000 \\
drm\_templated & e5-large & kmeans & importance\_weighted & 0 & 3200 & 4.00 & 0.965 & 0.999 & 1.000 \\
\end{longtable}

\begin{longtable}{lcccccc}
\caption{E3 per-cell results (learned baselines: PQ, OPQ, LSH, PCA+int8, HNSW-prune). (showing first 60 rows; full CSV in repo)}\label{tab:supp_e3}\\
\toprule
corpus & family & bytes\_per\_vec & identity\_accuracy & recall@1 & recall@10 & recall@100 \\
\midrule
\endfirsthead
\multicolumn{7}{l}{\small\itshape (continued)}\\
\toprule
corpus & family & bytes\_per\_vec & identity\_accuracy & recall@1 & recall@10 & recall@100 \\
\midrule
\endhead
\bottomrule
\endfoot
arxiv\_titles & hnsw\_prune & 4096 & 0.466 & 0.466 & 0.872 & 0.994 \\
arxiv\_titles & hnsw\_prune & 4096 & 0.518 & 0.517 & 0.877 & 0.995 \\
arxiv\_titles & hnsw\_prune & 4096 & 0.594 & 0.594 & 0.912 & 0.994 \\
arxiv\_titles & hnsw\_prune & 4096 & 0.668 & 0.670 & 0.941 & 0.996 \\
arxiv\_titles & lsh & 1 & 0.127 & 0.129 & 0.693 & 0.997 \\
arxiv\_titles & lsh & 4 & 0.310 & 0.311 & 0.834 & 0.998 \\
arxiv\_titles & lsh & 16 & 0.529 & 0.526 & 0.935 & 0.999 \\
arxiv\_titles & lsh & 64 & 0.684 & 0.686 & 0.953 & 0.999 \\
arxiv\_titles & opq & 4 & 0.679 & 0.679 & 0.938 & 0.995 \\
arxiv\_titles & opq & 8 & 0.699 & 0.699 & 0.945 & 0.996 \\
arxiv\_titles & opq & 16 & 0.709 & 0.710 & 0.956 & 0.996 \\
arxiv\_titles & opq & 32 & 0.710 & 0.710 & 0.955 & 0.996 \\
arxiv\_titles & opq & 64 & 0.716 & 0.720 & 0.954 & 0.997 \\
arxiv\_titles & pca\_int8 & 2 & 0.315 & 0.306 & 0.867 & 0.991 \\
arxiv\_titles & pca\_int8 & 4 & 0.450 & 0.448 & 0.911 & 0.993 \\
arxiv\_titles & pca\_int8 & 8 & 0.611 & 0.615 & 0.937 & 0.993 \\
arxiv\_titles & pca\_int8 & 16 & 0.682 & 0.683 & 0.946 & 0.995 \\
arxiv\_titles & pca\_int8 & 32 & 0.696 & 0.702 & 0.953 & 0.995 \\
arxiv\_titles & pca\_int8 & 64 & 0.718 & 0.717 & 0.951 & 0.995 \\
arxiv\_titles & pq & 4 & 0.664 & 0.663 & 0.911 & 0.991 \\
arxiv\_titles & pq & 8 & 0.692 & 0.687 & 0.917 & 0.990 \\
arxiv\_titles & pq & 16 & 0.710 & 0.709 & 0.930 & 0.991 \\
arxiv\_titles & pq & 32 & 0.716 & 0.720 & 0.937 & 0.991 \\
arxiv\_titles & pq & 64 & 0.731 & 0.731 & 0.950 & 0.994 \\
arxiv\_titles & pq & 128 & 0.720 & 0.723 & 0.956 & 0.997 \\
drm\_templated & hnsw\_prune & 4096 & 0.925 & 0.925 & 1.000 & 1.000 \\
drm\_templated & hnsw\_prune & 4096 & 0.930 & 0.930 & 1.000 & 1.000 \\
drm\_templated & hnsw\_prune & 4096 & 0.930 & 0.930 & 1.000 & 1.000 \\
drm\_templated & hnsw\_prune & 4096 & 0.930 & 0.930 & 1.000 & 1.000 \\
drm\_templated & lsh & 1 & 0.031 & 0.043 & 0.247 & 0.818 \\
drm\_templated & lsh & 4 & 0.514 & 0.512 & 0.874 & 0.999 \\
drm\_templated & lsh & 16 & 0.794 & 0.794 & 0.924 & 1.000 \\
drm\_templated & lsh & 64 & 0.805 & 0.805 & 0.944 & 1.000 \\
drm\_templated & opq & 4 & 0.710 & 0.705 & 0.866 & 1.000 \\
drm\_templated & opq & 8 & 0.746 & 0.746 & 0.906 & 1.000 \\
drm\_templated & opq & 16 & 0.769 & 0.769 & 0.927 & 1.000 \\
drm\_templated & opq & 32 & 0.780 & 0.780 & 0.929 & 1.000 \\
drm\_templated & opq & 64 & 0.797 & 0.797 & 0.934 & 1.000 \\
drm\_templated & pca\_int8 & 2 & 0.039 & 0.039 & 0.230 & 0.762 \\
drm\_templated & pca\_int8 & 4 & 0.250 & 0.250 & 0.655 & 0.995 \\
drm\_templated & pca\_int8 & 8 & 0.664 & 0.664 & 0.834 & 1.000 \\
drm\_templated & pca\_int8 & 16 & 0.731 & 0.731 & 0.875 & 1.000 \\
drm\_templated & pca\_int8 & 32 & 0.765 & 0.765 & 0.907 & 1.000 \\
drm\_templated & pca\_int8 & 64 & 0.799 & 0.799 & 0.931 & 1.000 \\
drm\_templated & pq & 4 & 0.517 & 0.510 & 0.767 & 1.000 \\
drm\_templated & pq & 8 & 0.675 & 0.662 & 0.804 & 1.000 \\
drm\_templated & pq & 16 & 0.738 & 0.738 & 0.824 & 1.000 \\
drm\_templated & pq & 32 & 0.755 & 0.755 & 0.843 & 1.000 \\
drm\_templated & pq & 64 & 0.772 & 0.772 & 0.871 & 1.000 \\
drm\_templated & pq & 128 & 0.794 & 0.794 & 0.919 & 1.000 \\
hotpot\_qa & hnsw\_prune & 4096 & 0.347 & 0.347 & 0.624 & 0.863 \\
hotpot\_qa & hnsw\_prune & 4096 & 0.347 & 0.347 & 0.624 & 0.863 \\
hotpot\_qa & hnsw\_prune & 4096 & 0.348 & 0.348 & 0.628 & 0.863 \\
hotpot\_qa & hnsw\_prune & 4096 & 0.366 & 0.365 & 0.649 & 0.880 \\
hotpot\_qa & lsh & 1 & 0.003 & 0.003 & 0.040 & 0.261 \\
hotpot\_qa & lsh & 4 & 0.032 & 0.032 & 0.146 & 0.462 \\
hotpot\_qa & lsh & 16 & 0.223 & 0.224 & 0.438 & 0.745 \\
hotpot\_qa & lsh & 64 & 0.372 & 0.373 & 0.655 & 0.894 \\
hotpot\_qa & opq & 4 & 0.205 & 0.206 & 0.491 & 0.848 \\
hotpot\_qa & opq & 8 & 0.276 & 0.277 & 0.574 & 0.874 \\
\end{longtable}

\begin{longtable}{lcccccccc}
\caption{E6 per-cell results (GAC router ablation). (showing first 60 rows; full CSV in repo)}\label{tab:supp_e6}\\
\toprule
corpus & model & router & seed & compression & identity\_accuracy & recall@10 & mrr@20 & coverage@0.80 \\
\midrule
\endfirsthead
\multicolumn{9}{l}{\small\itshape (continued)}\\
\toprule
corpus & model & router & seed & compression & identity\_accuracy & recall@10 & mrr@20 & coverage@0.80 \\
\midrule
\endhead
\bottomrule
\endfoot
arxiv\_titles & bge-large & gac\_fixed\_centroid & 0 & 2322.27 & 0.754 & 1.000 & 0.851 & 0.733 \\
arxiv\_titles & bge-large & gac\_fixed\_centroid & 1 & 2322.27 & 0.753 & 1.000 & 0.850 & 0.760 \\
arxiv\_titles & bge-large & gac\_fixed\_centroid & 2 & 2322.27 & 0.756 & 1.000 & 0.852 & 0.747 \\
arxiv\_titles & bge-large & gac\_fixed\_medoid & 0 & 580.57 & 0.624 & 0.908 & 0.703 & 0.122 \\
arxiv\_titles & bge-large & gac\_fixed\_medoid & 1 & 580.57 & 0.638 & 0.906 & 0.713 & 0.118 \\
arxiv\_titles & bge-large & gac\_fixed\_medoid & 2 & 580.57 & 0.627 & 0.906 & 0.703 & 0.113 \\
arxiv\_titles & bge-large & gac\_fixed\_prune & 0 & 2.00 & 0.582 & 0.938 & 0.707 & 0.612 \\
arxiv\_titles & bge-large & gac\_fixed\_prune & 1 & 2.00 & 0.583 & 0.936 & 0.709 & 0.630 \\
arxiv\_titles & bge-large & gac\_fixed\_prune & 2 & 2.00 & 0.587 & 0.943 & 0.712 & 0.611 \\
arxiv\_titles & bge-large & gac\_full & 0 & 2.00 & 0.582 & 0.938 & 0.707 & 0.612 \\
arxiv\_titles & bge-large & gac\_full & 1 & 2.00 & 0.583 & 0.936 & 0.709 & 0.630 \\
arxiv\_titles & bge-large & gac\_full & 2 & 2.00 & 0.587 & 0.943 & 0.712 & 0.611 \\
arxiv\_titles & bge-large & gac\_no\_residual & 0 & 2.00 & 0.582 & 0.938 & 0.707 & 0.612 \\
arxiv\_titles & bge-large & gac\_no\_residual & 1 & 2.00 & 0.583 & 0.936 & 0.709 & 0.630 \\
arxiv\_titles & bge-large & gac\_no\_residual & 2 & 2.00 & 0.587 & 0.943 & 0.712 & 0.611 \\
arxiv\_titles & bge-large & gac\_oracle & 0 & 2.38 & 0.594 & 0.949 & 0.721 & 0.689 \\
arxiv\_titles & bge-large & gac\_oracle & 1 & 2.38 & 0.611 & 0.955 & 0.733 & 0.707 \\
arxiv\_titles & bge-large & gac\_oracle & 2 & 2.16 & 0.588 & 0.947 & 0.716 & 0.655 \\
arxiv\_titles & bge-large & gac\_random & 0 & 4.39 & 0.593 & 0.888 & 0.694 & 0.689 \\
arxiv\_titles & bge-large & gac\_random & 1 & 5.69 & 0.571 & 0.945 & 0.690 & 0.662 \\
arxiv\_titles & bge-large & gac\_random & 2 & 5.81 & 0.477 & 0.824 & 0.594 & 0.590 \\
drm\_templated & bge-large & gac\_fixed\_centroid & 0 & 4.00 & 0.943 & 1.000 & 0.964 & 1.000 \\
drm\_templated & bge-large & gac\_fixed\_centroid & 1 & 4.00 & 0.945 & 1.000 & 0.967 & 1.000 \\
drm\_templated & bge-large & gac\_fixed\_centroid & 2 & 4.00 & 0.941 & 1.000 & 0.963 & 1.000 \\
drm\_templated & bge-large & gac\_fixed\_medoid & 0 & 1.00 & 0.918 & 0.988 & 0.939 & 1.000 \\
drm\_templated & bge-large & gac\_fixed\_medoid & 1 & 1.00 & 0.931 & 0.984 & 0.946 & 1.000 \\
drm\_templated & bge-large & gac\_fixed\_medoid & 2 & 1.00 & 0.911 & 0.983 & 0.931 & 1.000 \\
drm\_templated & bge-large & gac\_fixed\_prune & 0 & 2.00 & 0.769 & 0.925 & 0.822 & 1.000 \\
drm\_templated & bge-large & gac\_fixed\_prune & 1 & 2.00 & 0.791 & 0.930 & 0.840 & 1.000 \\
drm\_templated & bge-large & gac\_fixed\_prune & 2 & 2.00 & 0.768 & 0.924 & 0.823 & 1.000 \\
drm\_templated & bge-large & gac\_full & 0 & 1.55 & 0.932 & 0.995 & 0.953 & 1.000 \\
drm\_templated & bge-large & gac\_full & 1 & 1.53 & 0.946 & 0.995 & 0.964 & 1.000 \\
drm\_templated & bge-large & gac\_full & 2 & 1.54 & 0.943 & 0.994 & 0.960 & 1.000 \\
drm\_templated & bge-large & gac\_no\_residual & 0 & 4.00 & 0.934 & 0.995 & 0.956 & 1.000 \\
drm\_templated & bge-large & gac\_no\_residual & 1 & 4.00 & 0.949 & 0.996 & 0.966 & 1.000 \\
drm\_templated & bge-large & gac\_no\_residual & 2 & 4.00 & 0.944 & 0.996 & 0.962 & 1.000 \\
drm\_templated & bge-large & gac\_oracle & 0 & 4.00 & 0.943 & 1.000 & 0.964 & 1.000 \\
drm\_templated & bge-large & gac\_oracle & 1 & 4.00 & 0.945 & 1.000 & 0.967 & 1.000 \\
drm\_templated & bge-large & gac\_oracle & 2 & 4.00 & 0.941 & 1.000 & 0.963 & 1.000 \\
drm\_templated & bge-large & gac\_random & 0 & 1.77 & 0.821 & 0.969 & 0.868 & 1.000 \\
drm\_templated & bge-large & gac\_random & 1 & 1.73 & 0.866 & 0.978 & 0.898 & 1.000 \\
drm\_templated & bge-large & gac\_random & 2 & 1.74 & 0.834 & 0.969 & 0.875 & 1.000 \\
hotpot\_qa & bge-large & gac\_fixed\_centroid & 0 & 2.33 & 0.487 & 0.766 & 0.580 & 0.020 \\
hotpot\_qa & bge-large & gac\_fixed\_centroid & 1 & 2.33 & 0.518 & 0.797 & 0.610 & 0.023 \\
hotpot\_qa & bge-large & gac\_fixed\_centroid & 2 & 2.33 & 0.492 & 0.768 & 0.586 & 0.022 \\
hotpot\_qa & bge-large & gac\_fixed\_medoid & 0 & 1.01 & 0.361 & 0.572 & 0.419 & 0.025 \\
hotpot\_qa & bge-large & gac\_fixed\_medoid & 1 & 1.01 & 0.366 & 0.596 & 0.432 & 0.022 \\
hotpot\_qa & bge-large & gac\_fixed\_medoid & 2 & 1.01 & 0.345 & 0.576 & 0.406 & 0.018 \\
hotpot\_qa & bge-large & gac\_fixed\_prune & 0 & 1.81 & 0.352 & 0.623 & 0.439 & 0.018 \\
hotpot\_qa & bge-large & gac\_fixed\_prune & 1 & 1.81 & 0.366 & 0.632 & 0.447 & 0.018 \\
hotpot\_qa & bge-large & gac\_fixed\_prune & 2 & 1.81 & 0.348 & 0.620 & 0.434 & 0.020 \\
hotpot\_qa & bge-large & gac\_full & 0 & 1.79 & 0.355 & 0.628 & 0.440 & 0.018 \\
hotpot\_qa & bge-large & gac\_full & 1 & 1.79 & 0.360 & 0.634 & 0.450 & 0.018 \\
hotpot\_qa & bge-large & gac\_full & 2 & 1.80 & 0.354 & 0.633 & 0.442 & 0.019 \\
hotpot\_qa & bge-large & gac\_no\_residual & 0 & 1.81 & 0.351 & 0.624 & 0.439 & 0.018 \\
hotpot\_qa & bge-large & gac\_no\_residual & 1 & 1.81 & 0.367 & 0.633 & 0.447 & 0.018 \\
hotpot\_qa & bge-large & gac\_no\_residual & 2 & 1.81 & 0.348 & 0.622 & 0.435 & 0.021 \\
hotpot\_qa & bge-large & gac\_oracle & 0 & 2.11 & 0.471 & 0.744 & 0.564 & 0.021 \\
hotpot\_qa & bge-large & gac\_oracle & 1 & 2.13 & 0.510 & 0.770 & 0.598 & 0.024 \\
hotpot\_qa & bge-large & gac\_oracle & 2 & 2.15 & 0.487 & 0.749 & 0.576 & 0.022 \\
\end{longtable}


\section{DRM per-strategy breakdown}
\label{sec:supp:drm}

The DRM (Deese--Roediger--McDermott) corpus contains 64 templated
paraphrase lists of 50 items each --- the cleanest tight-regime case
in the paper, with cluster-local $\deff = 2.3$. Table~\ref{tab:drm_breakdown}
reports identity (literal and paraphrase), recall@1, and MRR per
strategy on bge-large. All four non-pruning strategies agree to within
$0.03$ identity on DRM; \textsc{selective\_prune} is the only
strategy that loses, consistent with the main paper's statement that
in the tight regime every reasonable strategy saturates.

\begin{table}[t]
\centering
\small
\caption{DRM templated: identity and MRR per strategy (bge-large, 3 seeds).}
\label{tab:drm_breakdown}
\begin{tabular}{lcccc}
\toprule
strategy & id\_lit & id\_para & recall1 & mrr \\
\midrule
centroid & 0.943 & 0.942 & 0.943 & 0.965 \\
gac & 0.940 & 0.938 & 0.940 & 0.959 \\
importance\_weighted & 0.941 & 0.941 & 0.941 & 0.963 \\
medoid & 0.911 & 0.909 & 0.911 & 0.944 \\
selective\_prune & 0.776 & 0.776 & 0.776 & 0.828 \\
\bottomrule
\end{tabular}
\end{table}


\section{Reproduction recipe}
\label{sec:supp:repro}

\paragraph{Local (single machine).}
\begin{verbatim}
git clone https://github.com/niashwin/geometry-of-consolidation
cd geometry-of-consolidation
pip install -r requirements.txt
export GAC_DATA_DIR=./data_local
python -m data.build_nq --model bge-large
python -m experiments.e1_theorem_validation --seed 0 --n-shards 1
python scripts/make_figures.py
\end{verbatim}

\paragraph{Modal (full reproduction).}
\begin{verbatim}
modal volume create gac-data
modal secret create gac-hf-secret HF_TOKEN=<your hf token>
modal run modal_app/app.py::run_all --run-id reproduce-v1
modal volume get gac-data runs/e1_theorem_validation/reproduce-v1 ./results/e1
\end{verbatim}

\paragraph{Seeds.}
The main paper uses seeds $\{0,1,2\}$ for experiments with three
replicates and $\{0,\dots,9\}$ for E1. Every parquet stores the seed
per row.

\section{Intentional scope choices}
\label{sec:supp:scope}

Three choices that a reviewer might reasonably ask about; we justify
each with the data that led to it rather than with a claim.

\paragraph{Embed-space paraphrase proxy (E4).}
E4 uses additive Gaussian perturbations in embedding space, calibrated
to three severity levels, as a paraphrase proxy rather than
LLM-generated paraphrases. On a preliminary DRM subset we compared
embed-space perturbation against Llama-3.1-70B paraphrase drift and
the two agreed within $5\%$ on identity error at matched severity.
Running 70B generation for $80{,}000$ paraphrases on Modal is a full
day of H100 time, which we prioritised for E5 (scale) and E7 (RAG)
instead. This is limitation L2 in the main paper.

\paragraph{Scale ceiling at 1M (E5).}
E5 sweeps $10{,}000 / 100{,}000 / 1{,}000{,}000$ Wikipedia items. We
stopped at 1M rather than $10{,}000{,}000$ because the identity
accuracy of GAC and centroid both flatten between $100{,}000$ and
$1{,}000{,}000$ on two encoders (BGE-large, MiniLM-L6); the
$1M \to 10M$ extrapolation from the observed curves gives no evidence
of a regime change. An additional 10M point would have required
$\approx 16$ H100-hours and a larger Modal volume; we did not spend
either.

\paragraph{No closed-weight encoder.}
E8 does not include \texttt{text-embedding-3-large} or other
closed-weight encoders. The six open-weight encoders already agree on
the strategy ordering and on the tight-regime $\eps_{\mathrm{cap}}
\leq 0.5\%$ result; adding one closed-weight encoder would not change
the qualitative conclusions and would add $\approx \$200$ of API spend
per full sweep.

\section{Run manifest}
\label{sec:supp:manifest}

Every figure and table in the main paper is produced from a single
parquet file generated on Modal. The manifest below records the
source script, parquet path, and commit SHA for each artifact so that
the pipeline can be audited end-to-end.

\begin{table}[t]
\centering
\small
\caption{Run manifest: every figure and table mapped to source parquet and commit SHA.}
\label{tab:manifest}
\begin{tabular}{lllc}
\toprule
artifact & source & parquet & commit \\
\midrule
Fig 1 (conceptual) & scripts/make\_figures.py:fig1b\_conceptual & (none) & 3b8a227477 \\
Fig 1 (theorem regimes) & scripts/make\_figures.py:fig1\_theorem\_regimes & results/e1/e1\_results.parquet & 3b8a227477 \\
Fig 2 & fig2\_strategy\_sweep & results/e4/e4\_results.parquet & 3b8a227477 \\
Fig 3 & fig3\_learned\_baselines & results/e3/e3\_results.parquet & 3b8a227477 \\
Fig 4a & fig4\_encoder\_universality & results/e8/e8\_results.parquet & 3b8a227477 \\
Fig 4b (id-coverage) & fig4b\_id\_coverage & results/e4/e4\_results.parquet & 3b8a227477 \\
Fig 5 (scale) & fig5\_scale & results/e5/e5\_results.parquet & 3b8a227477 \\
Fig 6 (ablation) & fig6\_ablation & results/e6/e6\_results.parquet & 3b8a227477 \\
Fig 6b (downstream) & fig6b\_downstream & results/e7/e7\_results.parquet & 3b8a227477 \\
Table 1 & make\_tables.py:t1 & e3, e4 parquets & 3b8a227477 \\
Table 2 & make\_tables.py:t2 & results/e6/e6\_results.parquet & 3b8a227477 \\
\bottomrule
\end{tabular}
\end{table}


\section{Licensing and data provenance}

All corpora used in the paper are public and licensed under MIT,
CC-BY-SA, or Apache-2.0. No human-subjects data is used. Per-corpus
URLs and licence terms are listed in the ``Data and code
availability'' section of the main paper.

\bibliographystyle{unsrtnat}
\bibliography{references}

\end{document}
